\chapter{Operazioni sul cluster}\label{chap:cluster}
\section{Setup online sul login node}
\begin{lstlisting}[language=bash]
ssh <user>@gcluster.dmi.unict.it
cd ~/neuro_symbolic_t2g
bash cluster/setup.sh
\end{lstlisting}
Il setup online esegue Apptainer \pathcode{/shared/sifs/latest.sif} direttamente
sul login node: niente \pathcode{srun}, allocazione GPU o Python host. Installa in
\pathcode{~/.local} senza alterare lo stack torch/CUDA, verifica versioni, scarica
Qwen e ASLG-PC12 e costruisce vocabolario/sidecar. \pathcode{BUILD_BIGRAM=1}
abilita la matrice opzionale.
Le versioni presidiate includono Transformers 5.3.0, TRL 0.24.0, PEFT 0.19.1,
Unsloth 2026.7.1 e torchao 0.16--0.17.

\section{Compute rigorosamente offline}
Preflight, train, eval e probe esportano prima di Python:
\begin{lstlisting}
HF_HUB_OFFLINE=1
TRANSFORMERS_OFFLINE=1
HF_DATASETS_OFFLINE=1
WANDB_MODE=offline
WANDB_DISABLE_WEAVE=true
WANDB_SILENT=true
\end{lstlisting}
Il preflight controlla import, snapshot, dataset, vocab, retrieval/bigram se
richiesti e W\&B. Non scarica e non allena.

\section{Coda e campaign chain}
La QoS ammette un solo job submitted/attivo per utente: \pathcode{run_all.sh}
mantiene una chain sequenziale di 12 voci. Non lanciare una seconda campagna in
parallelo. Comandi:
\begin{lstlisting}[language=bash]
CONFIG=experiments/configs/qwen25-05b/sft-grpo/zero-shot.yaml sbatch cluster/preflight.sh
bash cluster/run_all.sh
CONFIG=experiments/configs/qwen25-05b/sft-grpo/zero-shot.yaml EXTRA_ARGS="--resume" sbatch cluster/train.sh
\end{lstlisting}
Usare gli strumenti \pathcode{chain-show}/monitor per osservare l'avanzamento.
Una failure va corretta e ripresa con la stessa identità, non aggirata cambiando
path. PDA e reward usano i gate aggiuntivi descritti nei capitoli 6 e 7.
